\section{Tetrahedron: the largest-face bound}
The following elementary argument works for every nondegenerate tetrahedron,
not only a regular one.  For three points define the oriented area vector by
\[
 F(P,Q,R)=\frac12(Q-P)\mathbin\times(R-P).
\]
Its norm is the area of the triangle.  This vector is affine in each of
its three arguments separately, as is clear from
$2F(P,Q,R)=P\times Q+Q\times R+R\times P$.
Consequently, if $P,Q,R$ are convex combinations of the four vertices,
then $F(P,Q,R)$ is a convex combination of the $4^3=64$ corresponding
vertex-triangle vectors.  The norm of each of those vectors is at most
$M$, the largest face area.  The triangle inequality therefore proves
that every triangle contained in the tetrahedron has area at most $M$.

It remains to consider a quadrilateral section.  Label the vertices
$A,B$ on one side and $C,D$ on the other, with respective positive
signed heights $a,b,-c,-d$.  Translate so that $A=0$, and put
$B=u,C=v,D=w$.  The vertices of the section in their cyclic order are
\[
 P=\alpha v,\quad Q=\beta w,\quad
 R=(1-\gamma)u+\gamma w,\quad
 S=(1-\delta)u+\delta v,
\]
where
\[
 \alpha=\frac a{a+c},\qquad \beta=\frac a{a+d},\qquad
 \gamma=\frac b{b+d},\qquad \delta=\frac b{b+c}.
\]
Set
\[
 E_1=\frac12u\times v,\quad E_2=-\frac12u\times w,
 \quad E_3=\frac12v\times w,\quad E_4=-E_1-E_2-E_3.
\]
These are signed face-area vectors, so $\lVert E_i\rVert\le M$.
The quadrilateral area vector is
\[
 G=\frac12(R-P)\times(S-Q)=pE_1+qE_2+rE_3,
\]
with
\begin{align*}
 p&=\alpha+\delta-\alpha\delta-\gamma\delta,\\
 q&=\beta+\gamma-\beta\gamma-\gamma\delta,\\
 r&=\alpha\beta-\gamma\delta.
\end{align*}
The four numbers $p,q,p-r,q-r$ are nonnegative.  For example,
\[
 p=\alpha(1-\delta)+\delta(1-\gamma),\qquad
 p-r=\alpha(1-\beta)+\delta(1-\alpha),
\]
and the two analogous formulas give the claim for $q,q-r$.
If $r\le0$, apply the triangle inequality to $G$ as displayed.
If $r\ge0$, instead use
\[
 G=(p-r)E_1+(q-r)E_2-rE_4.
\]
In both cases,
\[
 \lVert G\rVert\le(p+q-r)M.
\]
Finally the exact slack identity is
\begin{align*}
 1-(p+q-r)
 &=(\alpha+\gamma-1)(\beta+\delta-1)\\
 &=\frac{(ab-cd)^2}{(a+c)(a+d)(b+c)(b+d)}\ge0.
\end{align*}
This proves the desired bound.  Sections with a vertex of the tetrahedron
on the cutting plane are triangles or degenerate quadrilaterals.  In
the latter case the same argument holds with nonnegative heights,
provided the four displayed denominators are positive; if a denominator
vanishes, the plane contains that entire edge, and the section is a
triangle.  Thus no limiting argument is needed.

For the regular model
\[
 T=\operatorname{conv}\{(1,1,1),(1,-1,-1),(-1,1,-1),(-1,-1,1)\},
\]
the edge length is $2\sqrt2$.  The plane $x+y+z=-1$ gives a face of
area $2\sqrt3$.  With edge length $s$ the maximum is
$\sqrt3\,s^2/4$.
Here a planar section includes the intersection with a supporting
plane.  If one requires the plane to meet the interior of $T$, this
number is only a supremum: every genuine triangular or quadrilateral
cut has strictly smaller area.  In the quadrilateral argument this
follows from the strict triangle inequality for two nonparallel,
nonzero face-area vectors; for a triangle the affine expansion gives
equality only for a complete face.

\section{Cube: an elementary central-section bound}
Let $C=[-1,1]^3$.  Assume the separately proved central-section
dominance lemma for centrally symmetric convex polytopes.  Sign changes
and permutations of coordinates reduce the normal to
$a\ge b\ge c\ge0$, with $a>0$.  Put $x=b/a$, $y=c/a$.
The orthogonal graph of the central section over the $yz$ square is
\[
 X=-xY-yZ,\qquad |Y|,|Z|\le1,\quad |xY+yZ|\le1.
\]
Area on this graph equals projected area times
$\sqrt{1+x^2+y^2}$.

If $x+y\le1$, there is no truncation, and the area is
\[
 A=4\sqrt{1+x^2+y^2}\le4\sqrt2,
\]
because $x^2+y^2\le(x+y)^2\le1$.

If $x+y>1$, let $r=x+y-1\in(0,1]$.  Here $x,y>0$.
Exactly two opposite corners are removed from the projected square;
each removed right triangle has legs $r/x,r/y$.  Thus
\[
 \frac A4=\sqrt{1+x^2+y^2}\left(1-\frac{r^2}{4xy}\right).
\]
The elementary inequalities
\[
 x^2+y^2\le1+r^2,\qquad 4xy\le(1+r)^2
\]
follow respectively from $(1-x)(1-y)\ge0$ and $(x-y)^2\ge0$.
Hence
\[
 \frac A4\le\sqrt{2+r^2}\frac{1+2r}{(1+r)^2}\le\sqrt2.
\]
The last step is justified by the explicit polynomial identity
\[
 2(1+r)^4-(2+r^2)(1+2r)^2
 =r^2(3+4r-2r^2)
 =r^2\bigl(3+2r+2r(1-r)\bigr)\ge0.
\]
All factors being squared or multiplied are nonnegative.
Equality is achieved at $(x,y)=(1,0)$, i.e. $X+Y=0$.
The section is the rectangle with vertices
$(1,-1,1),(1,-1,-1),(-1,1,-1),(-1,1,1)$,
with side lengths $2,2\sqrt2$ and area $4\sqrt2$.
At edge length $s$ the maximum is $\sqrt2\,s^2$.

\section{Octahedron: an all-positive polynomial certificate}
Let $O=\{(X,Y,Z):|X|+|Y|+|Z|\le1\}$, with edge length $\sqrt2$.
Again use central-section dominance.  Coordinate symmetries reduce a
nonzero normal to $(a,b,c)$ with $a,b,c\ge0$ and $a>0$.
If $b=c=0$, the section is the square $X=0$, of area $2$.
Otherwise $D=(a+b)(a+c)(b+c)>0$, and its vertices, allowing duplicates
when a coefficient vanishes, are
\[
 P=\frac{(b,-a,0)}{a+b},\qquad
 Q=\frac{(c,0,-a)}{a+c},\qquad
 R=\frac{(0,c,-b)}{b+c},\qquad -P,-Q,-R.
\]
They occur in the listed cyclic order.  The shoelace area vector is
\[
 P\times Q+Q\times R-R\times P
 =\frac{2(ab+ac+bc)}D(a,b,c).
\]
Consequently
\[
 A=\frac{2\sqrt{a^2+b^2+c^2}(ab+ac+bc)}D\le2,
\]
where the inequality follows by squaring and using
\begin{align*}
 D^2-(a^2+b^2+c^2)(ab+ac+bc)^2
 ={}&2(a^3b^3+a^3c^3+b^3c^3)\\
 &+4abc(a^2b+a^2c+ab^2+ac^2+b^2c+bc^2)\\
 &+7a^2b^2c^2\ge0.
\end{align*}
Every monomial has a nonnegative coefficient.  If at least two normal
coordinates are positive, the difference is strictly positive, so the
only maximizing central planes are the coordinate planes.
For example $Z=0$ gives the square with vertices
$(1,0,0),(0,1,0),(-1,0,0),(0,-1,0)$.
Scaling to edge length $s$ gives maximum area $s^2$.

\section{Arithmetic and formalization notes}
All inequalities above have fixed finite algebraic proofs over ordered
fields.  The displayed slack identities are checked by
\texttt{verify\_small\_solids.py}; the checker is supplementary evidence,
and the identities themselves are visible certificates.
For iPRA, integer and rational arithmetic are primitive recursive;
sign cases are decidable for encoded rationals or fixed-degree real
algebraic coordinates.  The operations and loops have fixed finite
bounds.  One must separately supply the representation and error
control for arbitrary constructive plane normals, as in the common
framework, rather than assume that arbitrary real-number sign tests
are decidable.  Square roots in the Euclidean area formulas can be
replaced throughout the upper-bound verification by nonnegative squared
areas.  In Lean, the polynomial identities are \texttt{ring}; their
nonnegativity uses \texttt{positivity} and \texttt{nlinarith} after the
explicit denominator-positivity lemmas.  The geometric bridge uses
finite convex combinations, cross products, and cyclic polygon area
vectors; it does not need an unbounded search or an existence theorem
for a maximizing plane.

The generic tetrahedron largest-face theorem is historically recorded
in M. S. Klamkin, ``Largest Cross-Section of a Tetrahedron'',
\emph{American Mathematical Monthly} 82 (1975), p.~661, solution to
Problem E~1298, pointing back to Problem~5006 (1962--1963).
An author-hosted collected reprint is at
\url{https://carmamaths.org/resources/jon/Preprints/Books/Klamkin/komplete.pdf}.
The proof given here is complete independently of that reference.
